% Apertura conceptual de la Parte I. Esta sección reúne, en una sola
% residencia, el problema histórico del continuo, la distinción entre marca,
% intervalo y medida, y la descomposición nonádica que conduce a APP.

\section{The Continuum Problem and the Genealogy of Measure}
\label{sec:continuo-numero-medida-hmt-revfinal}
\index{continuum!discrete structure}
\index{number!genealogy}
\index{measure!coupling}

This opening develops the number--measure dimension of the
\hyperref[sec:tesis-inaugural-helicoidal-hmt-md]{treatise's inaugural thesis}:
the discrete structure constructs the continuum and the real chart subsequently publishes
its coordinate, without intervening as a generating domain.

The history of the continuum can be read as the history of a separation and
of its successive reconciliations. Ancient arithmetic organized
pluralities; geometry compared lengths, areas, and volumes. The Greek
theory of proportions made it possible to reason with incommensurable magnitudes:
the diagonal of the unit square existed as an exact construction even when
its length could not be published as a ratio of integers. Books V and VII
of the \emph{Elements} therefore preserved two domains, magnitude and number,
united by rules of comparison but endowed with distinct operations
\cite{EuclidElements,SEPAristotleMath}.

Modern analysis brought that reconciliation to a decisive form. Cauchy sequences, Dedekind cuts, and set theory made \(\mathbb R\) a complete ordered field, capable of publishing continuous magnitudes in a single arithmetic language. Cantor distinguished cardinality from order; Dedekind made continuity a structural property of the number system; Weierstrass established the local discipline of limits and approximation \cite{SEPDedekind,SEPContinuity,SEPSetTheory}. The resulting real continuum is one of the most powerful constructions in the history of mathematics: it identifies presentations that converge to the same value and makes possible calculus, analytic geometry, and the quantitative formulation of physical theories.

Triadic Modular Holography poses a question prior to that
identification: \emph{what finite and prolongable structure produces the value
published by the real line?} Its primary object is the compatible genealogy of
an evaluation. A digit is the output of a positional reader; a value is
the scalar coordinate obtained by contracting a family of cylinders; an HMT
number is the section that additionally preserves phase, orientation, arithmetic
sheet, lifting quotient, carry, boundary, and route
memory. The Archimedean evaluation thus takes the form
\[
 \operatorname{ev}_{\mathbb R}:X_\infty^{\rm HMT}\longrightarrow\mathbb R,
\]
and their fibers gather stories that share a coordinate without thereby losing
their operational identity.

This formulation changes the explanatory order of measurement. Counting determines the cardinality of a finite set; measuring couples a state, a reader, and a chart, restricts their compatible continuations, and publishes an observable. If \(X\) is the state space, \(A\) the apparatus or reader, and \(\mathcal C\subset X\times A\) the compatibility condition, the minimal scheme is
\[
 (x,a)\in\mathcal C
 \longmapsto \mathcal O(x,a)
 \longmapsto \mathcal U\bigl(\mathcal O(x,a)\bigr),
\]
where \(\mathcal O\) conserves the type of the observable and \(\mathcal U\)
expresses it in the chosen chart. The final digit constitutes the visible coordinate
of the complete coupling. This distinction will be the mathematical basis of the
theory of the observer and of Dimensional Mechanics in Part V.

The epistemological consequence is precise. Ordinary physical metrology assigns coordinates to observables, and theories subsequently relate those coordinates by means of equations; HMT first constructs the state, the route, and the reader from which the coordinate descends. The fundamental constants then assume the status of results of a genealogy, and experimental comparison is applied to the output. In the same movement, an HMT symmetry expresses the conservation of unity through a transformation: the value records the published coordinate, and the exceptional projection records the organization that has survived transport. Structure and number become two verifiable aspects of the same discrete evolution.

\subsection{Marks, intervals, and the nonad cell}

The finite construction begins on the line of integers. The
origin \(0\), the decimal threshold \(10\), and the cell of positive
representatives are fixed
\[
 E_9=\{1,2,3,4,5,6,7,8,9\}.
\]
The elements of \(E_9\) perform two related but distinct functions: they are the nine phase marks of the path \(P_9\), and they are also the nine positive integer coordinates \(d(0,r)=r\) measured from the origin. This second reading adds no edges to the path. The nine marks determine exactly the eight half-open interior intervals
\[
 J_r=[r,r+1),\qquad 1\le r\le8.
\]
The closure edge is represented in the lifted chart by \([9,10)\); its nonadic projection returns to the first phase and preserves the carry of the next turn. The origin \(0\) belongs solely to the expanded decimal chart and does not add an interior interval to \(P_9\). Marks, coordinates, intervals, and closure are thus typed before any operation is applied.

For each \(n\ge1\), nonadic Euclidean division defines
\[
 \rho_9(n)=1+((n-1)\bmod9),
 \qquad
 q_9(n)=\left\lfloor\frac{n-1}{9}\right\rfloor,
\]
\begin{equation}
 n=9q_9(n)+\rho_9(n).
 \label{eq:rueda-helice-nonadica-revfinal}
\end{equation}
The coordinate \(\rho_9\) gives the position on a nine-phase wheel;
\(q_9\) preserves the accumulated height. The successor acts by
\[
 (q,r)\longmapsto
 \begin{cases}
 (q,r+1),&1\le r<9,\\
 (q+1,1),&r=9.
 \end{cases}
\]
Therefore, \(n\) and \(n+9\) share a phase and occupy consecutive heights. The
geometric realization
\[
 \mathcal H_9(n)=
 \left(
  \cos\frac{2\pi(\rho_9(n)-1)}9,
  \sin\frac{2\pi(\rho_9(n)-1)}9,
  q_9(n)+\frac{\rho_9(n)-1}{9}
 \right)
\]
makes this structure visible: projected onto the plane it is a wheel;
conserved in the three coordinates it is a helix. The return \(9\to1\)
recovers the angle and elevates the history. This is the elementary form of the
distinction that will govern the entire TPK:
\[
 \boxed{\text{phase return}\;\neq\;\text{state return}.}
\]

\subsection{From one cell to two coordinates: birth of APP}

The Cartesian product of two nonadic cells constructs a two-dimensional
chart. Its inner kernel contains \(8^2=64\) positions. The ninth
row and the ninth column publish the null residual class through the label
\(9\) and form a gnomon of seventeen cells:
\[
 81=64+(9+8)=64+17.
\]
The modular identification of opposite sides converts the square chart
into the torus
\[
 X_9=(\mathbb Z/9\mathbb Z)^2,
\]
and the cardinal translations endow it with a Cayley graph. The geometric
operation is, consequently, a precise sequence: kernel
\(8\times8\), gnomonic completion \(9\times9\), toroidal identification, and
transport by the graph.

On each vertex \((i,j)\in X_9\) act two evaluations of the same support:
the additive sheet \(\Sigma\) and the multiplicative sheet \(\Pi\). Their
exact incompatibility supplies dynamics because a route must preserve which
sheet acted, which residue it published and which quotient it produced. APP constructs the
support and its two evaluations; TRIT locally orients the transition; TPK
transports the complete state, simultaneously preserves both sheets and, through
\(M_{\rm ph}\), types the mode and TRIT phase of the transition; it also preserves the
memory of elevation and decides the survival of prolongations. The
inaugural architecture is fixed by the chain
\[
 \boxed{
 E_9\longrightarrow
 \mathrm{APP}\longrightarrow
 \mathrm{TRIT}\longrightarrow
 \mathrm{TPK}.}
\]

The importance of this beginning lies in its capacity for prolongation. The
same difference between position and height that turns the wheel into a helix is
reproduced in the periodicity of \(108=12\cdot9\) iterations, in the oriented
cyclic system of twelve phases, in the
\hyperref[sec:umbrales-prolongacion-tpk-inaugural]{architecture of levels and
thresholds of the TPK prolongation}
\(w_6\to w_{12}\to w_{18}\to w_{24}\to w_{30}\to R_{36}\to G_9\), in the derivation of
\(\pi,e,\varphi\) and \(\alpha_{\rm HMT}\), and in the double projection towards the
real value and exceptional incidence. The discrete structure of the continuum
designates precisely this continuity of a finite genealogy across
arbitrary depths.

\begin{statusbox}
The history of number, magnitude and continuum belongs to the classical mathematical
context. The decomposition \eqref{eq:rueda-helice-nonadica-revfinal}, the nonadic
torus, the APP sheets and memory transport are results of the
HMT corpus. The helical embedding \(\mathcal H_9\) is a geometric
formalization of that decomposition. Evaluation on all of \(\mathbb R\)
is established by the coverage and Archimedean closure theorem of the
genealogical continuum
\cref{thm:espiral-cierre-arquimediano-continuo-genealogico}; the genealogy of
the constructed values is formulated in the inverse system of compatible
sections.
\end{statusbox}
